\scanpage{65}
\mspage{56}
% archive sheet number: 65
pas clair que ce soit essentiel de le poser dans les définitions.

On veut d'autre part que $\unsure{f_M|f_\mu}$ soit compatible, en plus de $\leq$, à $\overset{\circ}{\ll},\mathrel{|\circ|}$. (\uline{NB} On sait déjà qu'il est compatible à $\unsure{\gtrless|\bowtie}$) Mais il faut lui donner un sens, dans le cas où \deleted{\unsure{on}} $f_M$ n'applique \inserted{pas} $\unsure{M|\mathcal{M}}$ dans $M'$, mais dans $\unsure{\mathscr{F}|\mathcal{F}|F}$ ! Mais il faut peut-être, avant cela, poser une autre condition, (comparer p. \unsure{24|28})

% vertical emphasis stroke beside the lower two lines of the condition
\[
\begin{aligned}
&\forall F\in\mathscr{F},\ \text{posant }F'=f(F),\ \text{l'application}\\
&\left|\begin{aligned}&G\longmapsto G'=f(F):\mathrm{Ss\,fig}(F)\overset{\sim}{\longrightarrow}\mathrm{Ss\,fig}\,F'\\
&\text{est un isom. d'ens ordonnés.}
\end{aligned}\right.\end{aligned}\tag*{(**)}
\]

\deleted{Utilisant le critère du \unsure{recollement|prolongement} des sous-figures, cela sera équivalent à :}
\[
\deleted{(**_M)\quad\forall X\in}
\]
Mais cela \unsure{indiquerait|impliquerait} que si $F$ est élémentaire, $F'$ aussi, c'est à dire finalement $f_M:M\longrightarrow M'$ ! C'est donc le moment de m'interroger si je connais des \deleted{\unsure{hom.}}\inserted{applications} $f:\mathscr{F}\longrightarrow\mathscr{F}'$ que j'aurais envie d'appeler ``homomorphismes'', et qui n'appliquent pas $M$ dans $M'$. Oui, il y a exemple \unsure{2|3} p. 50, mais du coup il m'apparaît hétéroclite etc.
\margin{\unsure{pas de quoi me laisser détourner de} (**).}
% margin note at the bottom edge, written horizontally
